% SOURCE_STATEMENT: "Problem 23. Among the many interesting open questions is the understanding of the nat-ural limits of the minimization method: (Alain Chenciner) - To what extent is it connected to symmetry constraints? - May topological constraints be imposed without forcing the minimizers to have collisions? - To what extent interesting results may be obtained for arbitrary masses? - What can be said of the hyperbolic and elliptic dimensions of a minimizer? - Is it interesting to look at minimizers with fixed energy?"
% Preparation and final self-audit: CONTRIBUTING.md prompts C and D; source/audit artifacts retained in the local physics_statements_2026_10_11 batch.
% !TeX program = lualatex
\documentclass[11pt]{article}
\usepackage[margin=25mm]{geometry}
\usepackage{fontspec}
\setmainfont{Latin Modern Roman}
\usepackage{amsmath,amssymb,mathtools}
\usepackage{unicode-math}
\setmathfont{Latin Modern Math}
\usepackage{xurl,hyperref}
\hypersetup{colorlinks=true,urlcolor=blue}
\setcounter{secnumdepth}{0}
\setlength{\parindent}{0pt}
\setlength{\parskip}{5pt}
\setlength{\emergencystretch}{3em}
\newcommand{\sref}[2]{\hyperlink{source:#1:#2}{[S#2]}}
\newcommand{\cataloguescope}[1]{\paragraph{Recorded target.}#1}
\newcommand{\R}{\mathbb R}
\newcommand{\C}{\mathbb C}
\newcommand{\N}{\mathbb N}
\newcommand{\Z}{\mathbb Z}
\newcommand{\Q}{\mathbb Q}
\newcommand{\F}{\mathbb F}
\newcommand{\eps}{\varepsilon}
\begin{document}
\section{Limits of action minimization: collisions, constraints, stability, and fixed energy}\label{statement}

\subsection{Definitions and mathematical statement}

\paragraph{Shared notation.}$\mathbb N=\{1,2,\ldots\}$; the standard number fields have their usual meanings. All additional conventions are stated locally.


\paragraph{Definitions and notation.}In the Newtonian $n$-body problem, positions $q_i\in\mathbb R^d$ and masses $m_i>0$ obey $m_i\ddot q_i=\nabla_{q_i}U$, where $U(q)=\sum_{i<j}m_im_j/|q_i-q_j|$ and the gravitational constant is normalized to one. Here $T=\tfrac12\sum_i m_i|\dot q_i|^2$, $E=T-U$, and a collision has $q_i=q_j$. The center of mass can be fixed at zero; moment of inertia is $I=\sum_i m_i|q_i|^2$. For a fixed period $P$, the action is $\mathcal A(q)=\int_0^P(T+U)\,dt$. A collision-free minimizer in an admissible loop class satisfies the Newtonian equations when variations are permitted. The period, constraints and closure of the class must be specified: a minimizing sequence can escape or approach collisions. Coercivity means action bounds control the chosen loop-space norm after removing allowed symmetries. Hyperbolic and elliptic directions of a periodic orbit refer to spectral stability of its linearized return map; the exact convention must address neutral symmetry directions. Fixed-energy variational principles differ from fixed-period action minimization.

\paragraph{Catalogue wording retained for traceability.}
Understand the natural limits of action minimization: determine its relationship to symmetry constraints; whether topological constraints can be imposed without forcing collisions; what results hold for arbitrary masses; the hyperbolic and elliptic dimensions of minimizers; and the usefulness of minimization at fixed energy.
\paragraph{Scope and source qualification.}All five programme questions are retained. A coercivity certificate in one symmetry class or a local stability statement does not settle the entire methodological target. The catalogue title/status was paired with a machine-generated research attempt; that attempt is not peer-reviewed evidence of settlement. The statement covers the original source questions rather than only the title’s subcase.
\subsection{Short English statement}
Determine the scope of celestial-mechanics minimization methods across symmetries, topology, masses, stability and fixed-energy constraints.
\subsection{Sources}
\begin{itemize}
\item[\hypertarget{source:statement:1}{S1}] ULAM.AI and the UnsolvedMath Contributors, UnsolvedMath, catalogue entry. Catalogue curation under CC BY 4.0; definitions and English exposition added locally.
\url{https://huggingface.co/datasets/ulamai/UnsolvedMath}.
\item[\hypertarget{source:statement:2}{S2}] American Institute of Mathematics, Variational Methods in Celestial Mechanics; notes by Kuo-Chang Chen, version 22 June 2003, source item 23.
\url{https://aimath.org/WWN/varcelest/varcelest.pdf}.
\end{itemize}
\label{end:statement}
\end{document}
